\begin{table*}[t]
\centering
\small
\begin{tabular}{llp{4.2cm}p{5.6cm}}
\toprule
Group & Axis & Values & Function \\
\midrule
Classifying & A: category & job, ecommerce, gov, support, $\ldots$ & Sector of the impersonated service. \\
Classifying & B: vector & phishing\_clone, credential\_harvest, dark\_patterns, reward\_trap, authority\_impersonation, conversational\_deception, prompt\_injection, fake\_trust\_signals & Social-engineering pattern (\S\ref{sec:vectors}). \\
Classifying & C: salience & blatant, subtle & Conspicuousness of the red flag. \\
Classifying & D: pii\_target & critical, high, medium & Sensitivity tier the attacker solicits. \\
\midrule
Factor & E: pressure & none, urgency, social\_proof, authority, scarcity & Coercive cue layered on the attack. \\
Factor & F: prompt\_injection & none, visible\_text, hidden\_dom, fake\_system\_msg & Modality of injected instruction targeting the agent. \\
Factor & G: interaction & static\_form, multi\_step\_form, single\_turn\_chat, multi\_turn\_chat & Interface through which PII is solicited. \\
Factor & H: multi\_site & single\_origin, multi\_origin\_same\_task, email\_entry\_plus\_site & Whether the attack spans web origins. \\
\bottomrule
\end{tabular}
\caption{Eight-axis taxonomy of Scammer4U environments. Axes A--D classify each environment; axes E--H are the factors used for paired-sibling ablation. Each sibling differs from its parent on exactly one axis.}
\label{tab:axes}
\end{table*}


\section{The \textsc{Scammer4U} Benchmark}
\label{sec:benchmark}

\subsection{Threat Model}
\label{sec:threat}

We model a web-browsing agent provisioned with a full user PII
profile, partitioned into four sensitivity tiers: \emph{critical}
(SSN, payment card with CVV, bank account, password, API key, 2FA
code), \emph{high} (email, phone, address), \emph{medium} (name,
employer, date of birth), and \emph{low} (residual fields). The
agent operates without human supervision, reflecting the
consumer-default deployment mode. The attacker controls one or more
web origins reachable from the agent's task and may use any combination
of visible content, hidden DOM instructions, simulated system messages,
and chat interactions to extract PII; out-of-scope surfaces (model
weights, TLS, DNS, host), the consumer-default no-confirmation
rationale, and the per-env \texttt{/api/captured} measurement
instrument are detailed in Appendix~\ref{app:threat-extended}. Our
primary question is whether any critical-tier value reaches an
attacker-controlled endpoint.

\subsection{Benchmark Design}
\label{sec:design}

\textsc{Scammer4U} contains 91 adversarial environments and 10
benign-twin baselines. Benign twins are built by removing the attack
surface from an adversarial environment while preserving the
legitimate task. They serve as a counterfactual baseline: because
twins produce 0\% critical-tier leakage, any leakage on adversarial
environments is attack-attributable rather than incidental
form-filling.

The 91 environments are organised on an 8-axis taxonomy
(Table~\ref{tab:axes}): four classifying axes (category, vector,
salience, PII target) and four factor axes used for paired
ablation (pressure, prompt-injection style, interaction modality,
multi-site flow). The eight attack vectors on axis B are each
mapped to MITRE ATT\&CK \citep{mitre_attack_citation}, OWASP LLM
Top 10 \citep{owasp_llm_citation}, and the ENISA AI threat
landscape \citep{enisa_citation} (Appendix~\ref{app:mitre}).

The central design feature is the paired-sibling construction: 60
of the 91 environments are siblings of a parent, differing on
exactly one axis. Any leakage difference between a parent and its
sibling can therefore be attributed to that single factor
\citep{Rubin1974EstimatingCE,Antonakis2010OnMC}, enabling causal
claims that cross-environment comparisons cannot support.

\subsection{Environment Construction and Fidelity}
\label{sec:construction}

Environments are self-hosted Flask applications generated from
human-written design briefs using an LLM authoring model distinct
from all four evaluated agents. The brief fixes every axis value,
the attack surface, and the visible red flags before code
generation; the authoring model plays only a clerical role. All 91
environments were manually reviewed to confirm axis correctness and
task completability (Appendix~\ref{app:authoring}).

A blind $20$-vs-$20$ review by $16$ reviewers rated interleaved
screenshots of \textsc{Scammer4U} environments and real captured
phishing on visual believability, copy quality, and
likelihood-of-fooling. Reviewers could not discriminate above
chance ($47.5\%$, $n = 640$, Fleiss $\kappa = 0.02$) and rated our
environments at or above real captures on every axis
(Appendix~\ref{app:fidelity}).

\subsection{Evaluation Setup}
\label{sec:setup}

\paragraph{Models and conditions.}
We evaluate GPT-5 mini, Claude Haiku 4.5, Gemini 3 Flash, and Llama
4 Scout in a shared Playwright-based harness
\citep{playwright,zhou2024webarenarealisticwebenvironment,Koh2024VisualWebArenaEM,10.5555/3692070.3694608,NEURIPS2023_5950bf29,he-etal-2024-webvoyager}
with identical observation format, action space, and prompt
template. Each model is tested under four conditions: \textbf{C0}
(no guidance), \textbf{C1} (generic privacy reminder), \textbf{C2}
(phishing-aware checklist), and \textbf{C3} (pre-submission
reflection). We run $n = 5$ seeds per (environment, model,
condition) cell, targeting $91 \cdot 4 \cdot 5 \cdot 4 + 10 \cdot 4 \cdot 5 = 7{,}480$
sessions total.

\paragraph{Metrics.}
The primary metric is \textbf{PLR$_{\text{crit}}$}, the rate at
which any critical-tier value reaches an attacker endpoint, measured
by exact string match against outbound POST bodies; profile values
are constructed to be globally unique, so paraphrased leaks are not
counted and PLR$_{\text{crit}}$ is a lower bound on adversarial
extraction (matcher precision/recall against a human-labelled subset
in Appendix~\ref{app:matcher}). Secondary metrics are a tier-weighted
leakage rate (\textbf{PLR$_{\text{weighted}}$}; weight vector and
PII-profile composition in Appendix~\ref{app:profile}), attack success rate
conditional on reaching the attack surface (\textbf{ASR}; we declare
trap URLs per environment and gate ASR on the agent's navigation to
them, disentangling non-leakage caused by detection from non-leakage
caused by failure to reach the trap), task completion rate
(\textbf{TCR}), a \textbf{Defended} flag for sessions that either
refuse or safely complete the task without reaching the attack
surface, and a detection flag (\textbf{DR}) indicating whether the
agent's reasoning identified the site as suspicious.

\paragraph{Detection rate judge.}
DR is operationalised via LLM-as-judge
\citep{GU2026101253,Zhu2023JudgeLMFL,Li2024FromGT}, asking at each
reasoning step whether the agent flagged the site as suspicious.
DR measures what the agent \emph{says}, not what it knows
internally. The primary judge is GPT-4o-mini and the secondary is
Llama 4 Scout; neither shares a model family with three of the four
evaluated agents, and the single within-family pair (GPT-4o-mini /
GPT-5 mini) is bounded by per-model stratified inter-judge $\kappa$
and a both-judges-agree sensitivity. Inter-judge agreement is
$\kappa = 0.39$--$0.44$ across all four models, with the within-family
cell in the same band as the cross-family cells
(Appendix~\ref{app:dr-sensitivity}). A keyword-match heuristic is
retained as an auxiliary signal but overstates the gap; the LLM-judge
primary is the more conservative reading.

\paragraph{Analysis.}
The pre-registered primary family is a within-condition
detection--action test (F1), ten paired-sibling factor tests
(F2--F11), and three mitigation contrasts (M1--M3). Each is a
mixed-effects logistic regression on session-level
PLR$_{\text{crit}}$ with environment and model as crossed random
effects, with a Wilcoxon signed-rank fallback on convergence failure;
Benjamini--Hochberg correction at $q = 0.05$ is applied across the
inference-bearing subset selected by a pre-registered power gate
(Appendix~\ref{app:claims}). Point estimates are reported in
percentage points with 95\% CIs. Two claims carry pre-committed
falsification thresholds: the mitigation-insufficiency claim is
falsified if any condition reduces pooled PLR$_{\text{crit}}$ by
$\geq 30$\,pp across all four models; the detection--action claim is
falsified if within-C3 PLR$_{\text{crit}} \mid \text{DR}{=}1
\leq 0.10$ \citep{Barocas2021DesigningDE,MLSYS2021_0184b0cd}.

\paragraph{Pre-registration.}
The full analysis plan was committed and tagged
\texttt{prereg-v2-start} before data collection
\citep{Nosek2018ThePR,Nosek2019PreregistrationIH,Gallitto2024ExternalVO}.
An identified URL-shape contamination in an earlier harness (visible
localhost cues acting as model-side trust shortcuts) was removed
before the reported runs; the pre-fix slice is preserved for
contamination accounting (Appendix~\ref{app:contamination}), and the
correction alters no experimental-design choice. Post-tag deviations
are enumerated with their pre-fix vs post-fix slices in
Appendix~\ref{app:deviations}.